%% toc: P2 10 10.3,10.4
\section{Proofs for equivalent substitutions}
\label{sec:annex-p02-w02-p02-004-equivalent-substitutions}

Throughout, \(\operatorname{Var}(\varphi)\) and equivalence of substitutions on
a set of variables are as in
Definition~\ref{def:p02-w02-p02-004-occurring-variables}, and instances
\(\varphi[\upsilon]\) are as in
Definition~\ref{def:p02-uniform-substitution}.  The facts about \(\equiv\)
used below are reflexivity, transitivity and the congruence clauses for
\(\neg,\land,\lor,\to\), all proved in
Annex~\ref{annex:p2.prop.replacement}, and the last consequence of the
substitution theorem, proved in Annex~\ref{annex:p2.prop.substitution}.

\begin{annexitem}{annex:p02-w02-p02-004-eqsub}{thm:p02-w02-p02-004-equivalent-substitutions}{Equivalent substitutions, equivalent instances}
\begin{proof}
We prove, by structural induction on \(\varphi\), the statement: for all
uniform substitutions \(\upsilon,\upsilon'\) equivalent on
\(\operatorname{Var}(\varphi)\),
\(\varphi[\upsilon]\equiv\varphi[\upsilon']\).  The pair of substitutions is
universally quantified inside the induction statement, so each induction
hypothesis may be applied to the same pair.
Theorem~\ref{thm:prop-unique-readability} makes the case distinction below
exhaustive and exclusive.

\emph{Variable.}  Let \(\varphi=p_n\).  Then
\(\operatorname{Var}(\varphi)=\{p_n\}\), so the hypothesis gives
\(\upsilon(p_n)\equiv\upsilon'(p_n)\).  By the variable clause of
Definition~\ref{def:p02-uniform-substitution},
\(p_n[\upsilon]=\upsilon(p_n)\) and \(p_n[\upsilon']=\upsilon'(p_n)\), which
proves the claim.

\emph{Falsum.}  Both instances equal \(\bottom\), and
\(\bottom\equiv\bottom\) by reflexivity.  The hypothesis on the
substitutions is not used, because \(\operatorname{Var}(\bottom)=\varnothing\).

\emph{Negation.}  Let \(\varphi=\neg\alpha\).  Since
\(\operatorname{Var}(\alpha)=\operatorname{Var}(\neg\alpha)\), the
substitutions are equivalent on \(\operatorname{Var}(\alpha)\), and the
induction hypothesis gives \(\alpha[\upsilon]\equiv\alpha[\upsilon']\).
Negation congruence gives
\(\neg(\alpha[\upsilon])\equiv\neg(\alpha[\upsilon'])\).  By the negation
clause of the instance map these formulas are
\((\neg\alpha)[\upsilon]\) and \((\neg\alpha)[\upsilon']\).

\emph{Binary constructors.}  Let \(\varphi=(\alpha\circ\beta)\) with
\(\circ\in\{\land,\lor,\to\}\).  Both \(\operatorname{Var}(\alpha)\) and
\(\operatorname{Var}(\beta)\) are subsets of
\(\operatorname{Var}(\alpha\circ\beta)\), so the substitutions are
equivalent on each of them.  The two induction hypotheses give
\(\alpha[\upsilon]\equiv\alpha[\upsilon']\) and
\(\beta[\upsilon]\equiv\beta[\upsilon']\).  The congruence clause for
\(\circ\) gives
\[
 (\alpha[\upsilon]\circ\beta[\upsilon])
   \equiv(\alpha[\upsilon']\circ\beta[\upsilon']).
\]
By the clause of the instance map for \(\circ\), the two sides are
\((\alpha\circ\beta)[\upsilon]\) and \((\alpha\circ\beta)[\upsilon']\).
This is done for \(\land\), \(\lor\) and \(\to\) separately, each with its
own congruence clause.

These cases exhaust the primitive grammar.  No property of
\(\upsilon(p)\) or \(\upsilon'(p)\) for \(p\notin\operatorname{Var}(\varphi)\)
was used.
\end{proof}
\end{annexitem}

\begin{annexitem}{annex:p02-w02-p02-004-invar}{thm:p02-w02-p02-004-substitution-invariance}{Joint substitution invariance of equivalence}
\begin{proof}
Assume \(\varphi\equiv\varphi'\) and that \(\upsilon,\upsilon'\) are
equivalent on \(\operatorname{Var}(\varphi)\cup\operatorname{Var}(\varphi')\).

First, Theorem~\ref{thm:p02-uniform-substitution}, applied to the single
substitution \(\upsilon\) and the equivalent formulas \(\varphi,\varphi'\),
gives \(\varphi[\upsilon]\equiv\varphi'[\upsilon]\).

Second, \(\upsilon\) and \(\upsilon'\) are equivalent on the subset
\(\operatorname{Var}(\varphi')\), so
Theorem~\ref{thm:p02-w02-p02-004-equivalent-substitutions}, applied to
\(\varphi'\), gives \(\varphi'[\upsilon]\equiv\varphi'[\upsilon']\).

Transitivity of \(\equiv\) combines the two steps into
\(\varphi[\upsilon]\equiv\varphi'[\upsilon']\).

For the validity clause, assume \(\entails\varphi\) and that
\(\upsilon,\upsilon'\) are equivalent on \(\operatorname{Var}(\varphi)\).
The validity consequence of Theorem~\ref{thm:p02-uniform-substitution}
gives \(\entails\varphi[\upsilon]\), so every valuation \(v\) satisfies
\(\widehat v(\varphi[\upsilon])=\mathbf T\).
Theorem~\ref{thm:p02-w02-p02-004-equivalent-substitutions} gives
\(\varphi[\upsilon]\equiv\varphi[\upsilon']\), so for every valuation \(v\)
also \(\widehat v(\varphi[\upsilon'])=\widehat v(\varphi[\upsilon])
=\mathbf T\).  Hence \(\entails\varphi[\upsilon']\).
\end{proof}
\end{annexitem}
